\section*{第 \ref{ch:b1:electronic-effects} 章——电子效应与活性中间体}

\begin{solution}{exo:b1:electronic-effects:1}
2-甲基丁烷，\ce{CH3-CH(CH3)-CH2-CH3}：三个 \ce{CH3} 为伯碳，
\ce{CH2} 为仲碳，CH 为叔碳。2-溴-2-甲基丙烷中，三个
\ce{CH3} 为伯碳，中心碳为叔碳。四种溴化物
\ce{C4H9Br}（1-溴丁烷、2-溴丁烷、1-溴-2-甲基丙烷、
2-溴-2-甲基丙烷），其中最后一种生成叔碳正离子，
最稳定。
\end{solution}

\begin{solution}{exo:b1:electronic-effects:2}
$10^{4.76 - 2.87} = 10^{1.89} = 78$; $10^{4.76 - 0.52} = 10^{4.24} =
\num{1.7e4}$.
\end{solution}

\begin{solution}{exo:b1:electronic-effects:3}
乙酸根：C=O 与 \ce{C-O^-} 在两个氧原子之间互换。
4-硝基苯氧负离子：氧的孤电子对形成 C=O，环上的双键
移动，连接 \ce{NO2} 的碳形成 C=N 键，一对 N=O 电子
转移到氧上：得到一种结构，环的一端为 C=O，另一端为 C=N，
两个硝基氧均带负电（醌式结构）。烯丙基正离子：
\ce{CH2=CH-CH2+} $\leftrightarrow$ \ce{+CH2-CH=CH2}.
\end{solution}

\begin{solution}{exo:b1:electronic-effects:4}
第二个箭头：从水的 O--H 键指向其氧原子；生成物 \ce{H2O}
（由 \ce{OH-} 转变而来，现已呈中性）和 \ce{OH-}（原水分子的氧
保留电子对，电荷 $-1$）。对于氨：一个箭头从 N 的孤电子对指向
\ce{H3O+} 中的一个 H，另一个箭头从该 O--H 键指向氧：N 变为
$+1$，O 变为 0。
\end{solution}

\begin{solution}{exo:b1:electronic-effects:5}
乙醇（15.9）$<$ 水（14.0）$<$ 苯酚（9.99，环上离域）$<$
4-硝基苯酚（7.16，\ce{NO2} 承担电荷）$<$ 乙酸（4.76，
两个等价氧）$<$ 三氯乙酸（0.51，三个具有 $-I$
效应的氯原子）。
\end{solution}

\begin{solution}{exo:b1:electronic-effects:6}
苯胺（4.60）$<$ 吡啶（5.23）$<$ 氨（9.25）$<$ 甲胺
（10.66）。苯胺中氮上的孤电子对与苯环共享
（含 C=N$^+$ 的结构，负电荷位于 \textit{邻位}和
\textit{对位}）；使 N 质子化将失去这种离域。
\end{solution}

\begin{solution}{exo:b1:electronic-effects:7}
在 3-硝基苯氧负离子中，共振结构把电荷放在
与氧呈 \textit{邻位}和 \textit{对位}的碳原子上，从不在
连有 \ce{NO2} 的碳上，无共轭效应助力，因此弱于 4-异构体。
硝基仍通过 $\sigma$ 键吸引电子（$-I$），
因而强于苯酚。
\end{solution}

\begin{solution}{exo:b1:electronic-effects:8}
\ce{CH3+} $<$ \ce{CH3CH2+} $<$ \ce{CH2=CH-CH2+} $<$ \ce{C6H5CH2+}
$\approx$ \ce{(CH3)3C+} $<$ \ce{CH3OCH2+}：烷基给电子；
烯丙基和苄基使电荷离域（分别对应两种和四种共振结构）；
氧的孤电子对给出一种各原子都满足八电子的结构。
列表中间部分的顺序取决于介质。
\end{solution}

\begin{solution}{exo:b1:electronic-effects:9}
\ce{OH-}（亲核试剂）进攻 \ce{CH3Br} 中的碳（亲电试剂），
C--Br 键的电子对转移到 Br 上并离去。\ce{H2O} 向 \ce{H+} 提供一对电子（此处
酸与碱仅为所用的名称）。\ce{CN-} 进攻
乙醛的羰基碳，C=O 的 $\pi$ 电子对移到氧上。\ce{NH3} 将其孤电子对
提供给 \ce{BF3} 中硼的空轨道。
\end{solution}

\begin{solution}{exo:b1:electronic-effects:10}
否：$K_a = 10^{-0.51} = 0.31$ 远大于 $C$；该酸几乎
完全解离。由 $x^2/(0.010 - x) = 0.31$ 得 $x = \qty{9.7e-3}{mol/L}$，
$\mathrm{pH} = 2.01$，解离度为 \qty{97}{\percent}。
\end{solution}

\begin{solution}{exo:b1:electronic-effects:11}
烷基将电子（$+I$）推向醇盐的氧，使其
负电荷更不稳定，基团越多，这种作用越明显。较大的醇盐
也较难被水溶剂化；水对较小带电氧的稳定作用
强于埋在烷基之中的氧。两者都使
$\mathrm{p}K_a$ 升高。
\end{solution}

\begin{solution}{exo:b1:electronic-effects:12}
在酰胺中，氮的孤电子对与 C=O 共享
（\ce{N+=C-O^-} 结构）；若用它去结合质子或亲电试剂，
将破坏这种稳定作用。在强酸中，酰胺在
带部分负电荷的氧上被质子化。
\end{solution}

\begin{solution}{pb:b1:electronic-effects:1}
\textbf{1.} $10^{1.89} = 78$; $10^{1.52} = 33$; $10^{0.84} = 6.9$.
\textbf{2.} 否：每多一个氯，作用都更小。羧酸根的电荷
已充分分散，且随酸接近完全解离，测量
精度下降。
\textbf{3.} 每个 Cl 通过 C--C 和 C--O
键吸引电子密度，使羧酸根稳定（$-I$）。
\textbf{4.} 诱导效应在两三个键之内就会减弱。
\textbf{5.} 两种酸看起来强度相当（0.52 与 0.51）。在水中，
在通常浓度下，两者几乎都完全解离
（见\cref{ch:b1:predominant-reaction}，拉平效应）：$\mathrm{p}K_a$ 接近 0
是水能够区分的极限，两个数值
的不确定度约为零点几。
\textbf{6.} 乙酸：以酸式存在（\qty{98}{\percent}）；氯乙酸：
两种形式并存，阴离子略占优势（$10^{3 - 2.87} = 1.3$）；其余三种
均为阴离子形式。
\textbf{7.} \ce{CH3-C(=O)-O^-} $\leftrightarrow$ \ce{CH3-C(-O^-)=O}.
\textbf{8.} 两个等长的 C--O 键，介于双键和单键之间。
\textbf{9.} 电荷位于单键相连的氧上，该氧上没有
可与之共享的 $\pi$ 键。
\textbf{10.} $10^{15.9 - 4.76} = 10^{11.1}$，约为 $10^{11}$。
\textbf{11.} 苯氧负离子的电荷分散到三个环碳上，这些碳的
电负性小于氧，这是一种稳定化收益，小于羧酸根中的收益：
$\mathrm{p}K_a$ 9.99 介于 4.76 与 15.9 之间。
\textbf{12.} $K = 10^{6.37 - \mathrm{p}K_a}$：乙酸 $10^{1.61} = 41$，
反应有利（过量碳酸氢根时放出二氧化碳）；
苯酚 $10^{-3.62}$ 与乙醇 $10^{-9.5}$：不发生反应。
\textbf{13.} 4-硝基苯酚（7.16）$>$ 2-硝基苯酚（7.23）$>$ 3-硝基苯酚
（8.36）$>$ 苯酚（9.99）。
\textbf{14.} \cref{exo:b1:electronic-effects:3} 中的醌式结构，
电荷位于硝基的一个氧原子上。
\textbf{15.} 电荷只到达与氧呈 \textit{邻位}和
\textit{对位}的碳上；连硝基的碳为 \textit{间位}。
\textbf{16.} \ce{NO2} 的 $-I$ 效应。
\textbf{17.} 在 2-硝基苯酚中，羟基与相邻硝基的一个氧
形成氢键，它稳定酸式，并使
质子稍难被除去。
\textbf{18.} 阴离子分数为 $1/(1 + 10^{\mathrm{p}K_a - 7.5})$：4-硝基
\qty{69}{\percent}，2-硝基 \qty{65}{\percent}，3-硝基
\qty{12}{\percent}。4-硝基苯酚的酸式与阴离子吸收不同，
阴离子的吸收位于可见光区，其颜色在 $\mathrm{p}K_a$ 附近出现。
\textbf{19.} $x^2/(0.010 - x) = 10^{-0.52} = 0.30$：$x = \qty{9.7e-3}{mol/L}$，
$\mathrm{pH} = 2.01$：几乎是强酸。
\textbf{20.} $\frac12(4.76 + 2.00) = 3.38$.
\textbf{21.} \ce{CF3COOH + CH3COO- -> CF3COO- + CH3COOH}, $K = 10^{4.76 -
0.52} = 10^{4.24}$：定量反应。
\textbf{22.} $K_a(\ce{CF3COOH})/K_a(\ce{CH3COOH}) =$ \textbf{$10^{4.24} =
\num{1.7e4}$}，约一万倍。
\end{solution}
